\section{結言}
今回の話題を結ぶのは、AIの成果を人間や既存の技術体系の中でどう扱うかという問題である。数学の形式検証では証明の確認手続きが焦点となり、エージェントでは承認ゲート、情報アクセス権限、長時間待機の実装が議論された。モデル性能だけでなく、検証可能性と運用上の制御が価値を左右する流れが見える。

一方、モデル採用やベンチマーク、金融市場を巡る投稿には一次情報の不足が目立った。Grokによる今回の観測はPartialであり、窓外の発表を受けた二次的な議論も多い。翌日以降は公式仕様や独立テストの公開状況が注目点となる。
